% ===================================================================
% Paste these sections into the main .tex of the provided template.
% Everything marked TODO must be filled from results/results_table.md
% after running the Colab notebook. Do not change template formatting.
% ===================================================================

% ---- header info (put where the template asks for it) ----
% Team XX: Name 1 (Roll), Name 2 (Roll), Name 3 (Roll), Name 4 (Roll)
% GitHub: \url{https://github.com/<team>/ffs-diverse}

\section{Summary of the Assigned Paper}
Video-language models are costly because every frame goes through a large model, so systems sample a few frames
uniformly and can miss the moment that matters. FFS~\cite{buch2025ffs} puts a small selector over cheap CLIP features in
front of the large model. It picks question-relevant frames, and learnable \texttt{[CTRL]} tokens inside a
differentiable top-$k$ let it decide \emph{how many} ($k{-}c$ to $k$) each question needs. It learns from the answer loss
alone, with no frame labels, and keeps dense-frame accuracy at a fraction of the compute.

\section{Proposed Novelty: Redundancy-Aware Selection (RAS)}
FFS scores frames independently, so a set's value is additive, $\sum_{i\in A} l_i$. Neighbouring frames are
near-duplicates with near-equal scores, so top-$k$ spends its tiny budget on repeats, and \texttt{[CTRL]} judges a frame
by relevance, not by what it \emph{adds}. We replace top-$k$ by greedy selection with a learned redundancy penalty, using
a video-adaptive similarity of frozen CLIP features $f_i$:
\begin{equation}
S_{ij}=\tfrac{[\cos(f_i,f_j)-\mu]_+}{1-\mu},\qquad
i_t=\arg\max_{i\notin A_{t-1}}\big(\underbrace{l_i+\gamma_i}_{\text{FFS}}-\beta\textstyle\sum_{j\in A_{t-1}}S_{ij}\big),\;t=1..k,
\end{equation}
where $\mu$ is the video's mean off-diagonal cosine, $S{=}0$ for \texttt{[CTRL]}, $A_0=\varnothing$,
and $\gamma_i$ is Gumbel noise (training only); picked \texttt{[CTRL]}s are dropped as in FFS. Gradients use the relaxed subset
sampler~\cite{xie2019subset} with the same penalty, $w\leftarrow w+\log(1-p)-\tfrac{\beta}{\tau}Sp$, $p=\mathrm{softmax}(w)$,
plus a straight-through estimator; $\beta=\mathrm{softplus}(\rho)$ is learned. The \emph{only} difference from our FFS
re-implementation is $\beta$ ($\beta=0$ gives FFS exactly). Both share one recipe: count supervision from fixed-$k$ policies
(a \texttt{[CTRL]} hinge to $k^\star$, the fewest frames that answer correctly) and task gradients restricted to frames.

\section{Theoretical Justification}
\textbf{Diminishing returns.} Eq.~(1) is the greedy step for
$F(A)=\sum_{i\in A}l_i-\beta\sum_{\{i,j\}\subseteq A}S_{ij}$, with marginal gain
$\Delta(i\,|\,A)=l_i-\beta\sum_{j\in A}S_{ij}$. Because $S\ge 0$, $A\subseteq B\Rightarrow\Delta(i|A)\ge\Delta(i|B)$, so $F$ is
submodular. $F$ also has the shape of a DPP
objective~\cite{kulesza2012dpp}: with quality $q_i=e^{l_i}$ and unit-diagonal kernel $K$,
$\log\det L_A=2\sum_{A}l_i+\log\det K_A\approx 2\sum_{A}l_i-\sum_{\{i,j\}\subseteq A}K_{ij}^2$ (second order). That is
relevance minus pairwise redundancy, and RAS is its greedy (MMR-style) analogue with $S_{ij}$ in place of $K_{ij}^2$.
So it picks frames that are relevant \emph{and} diverse.
\textbf{Adaptive budget.} A frame is admitted at step $t$ only if $\max_i\Delta(i|A_{t-1})$ beats the next unused
\texttt{[CTRL]} score $c_{(r)}$. FFS compares raw relevance $l_i$ with $c_{(r)}$, so it counts \emph{relevant} frames.
RAS counts \emph{relevant and new} frames. For the same scores, $s_{\text{RAS}}\le s_{\text{FFS}}$: each of the
$m$ RAS frames had $l_i\ge\Delta\ge c_{(k-m+1)}$, so FFS admits at least $m$. They are equal e.g.\ when the picks are
mutually dissimilar. The saving grows with video redundancy, where uniform sampling and FFS waste compute.
\textbf{Safe extension.} $\beta{=}0$ is FFS and $\beta$ is learned, so training can switch the penalty off if useless.

\section{Comparative Results}
NExT-QA val (4{,}996 Qs), EVA02-CLIP-L/14 features, 32-frame pool, one shared frozen downstream model, 3 seeds.
\begin{center}\footnotesize
% PLACEHOLDER - replace with results/results_table_compact.tex produced by evaluate.py on Colab
\begin{tabular}{@{}lccccccc@{}}
\toprule
Method & Frames & Acc & C & T & D & Red.$\downarrow$ & Cost \\
\midrule
Uniform@4 & 4.00 & -- & -- & -- & -- & -- & 0.125 \\
Uniform@32 & 32.00 & -- & -- & -- & -- & -- & 1.000 \\
Fixed top-4 & 4.00 & -- & -- & -- & -- & -- & 0.125 \\
FFS (re-impl.) & -- & -- & -- & -- & -- & -- & -- \\
\textbf{FFS+RAS (ours)} & -- & -- & -- & -- & -- & -- & -- \\
\bottomrule
\end{tabular}
   % auto-generated by evaluate.py (results/results_table_compact.tex)
\\[2pt] Tab.\,1: NExT-QA val accuracy (\%) by question type (C/T/D), avg.\ frames given to the downstream model,
redundancy of picks, and downstream cost relative to 32 frames. Learned $\beta$ = TODO.
\end{center}
TODO: 2--3 sentences of interpretation (template below). Wider coverage should help multi-event temporal/causal questions;
fine-motion questions that need consecutive frames may lose.
% TODO (2-3 sentences, after running): e.g.
% "RAS improves accuracy over FFS by X points while using Y fewer frames on average (Z x less downstream compute than
%  uniform@32). Gains concentrate on Temporal/Causal questions; Descriptive questions change by W. Selected-frame
%  redundancy drops from A to B, confirming the mechanism. Trade-off: ..."

% ---- bib entries (add to the template's .bib or thebibliography) ----
% buch2025ffs: S. Buch, A. Nagrani, A. Arnab, C. Schmid. Flexible Frame Selection for Efficient Video Reasoning. CVPR 2025.
% xie2019subset: S. M. Xie, S. Ermon. Reparameterizable Subset Sampling via Continuous Relaxations. IJCAI 2019.
% kulesza2012dpp: A. Kulesza, B. Taskar. Determinantal Point Processes for Machine Learning. FnT ML 2012.
